\documentclass{harryopo-paper}

\title{边缘计算环境下基于Flink的实时数据流处理与动态资源调度优化}
\author{李四人工智能学院2025000102\\ \fzfs\small 示例大学}
\date{2026年09月02日}


\begin{document}

\maketitle
\begin{abstract}
随着物联网设备的快速普及和海量实时数据的产生，传统云计算架构在处理边缘侧实时数据时面临高延迟、带宽瓶颈和隐私保护不足等挑战。本文针对边缘计算环境下的实时数据流处理问题，提出了一种基于Apache Flink的分布式流处理框架，并设计了动态资源调度优化策略。实验结果表明，所提框架在毫秒级延迟约束下，吞吐量较基准方法提升了32\%，资源利用率提高了28\%。
\end{abstract}
\keywords{边缘计算；实时流处理；Apache Flink；动态资源调度}

\subsection*{一、引言}
\addcontentsline{toc}{subsection}{一、引言}

边缘计算作为一种新型计算范式，通过将计算能力从云端下沉到网络边缘，为终端设备提供就近服务。随着5G技术的商用化推进，边缘计算在智能制造、自动驾驶、智慧医疗等领域的应用前景愈发广阔。

实时数据流处理是边缘计算的核心需求之一。与传统批处理不同，流处理要求系统在毫秒级延迟内完成数据采集、过滤、聚合和响应。Apache Flink作为一个真正的流处理引擎，以其精确一次（Exactly-once）语义和低延迟特性，成为边缘计算场景下的首选流处理框架。

\subsection*{二、系统架构}
\addcontentsline{toc}{subsection}{二、系统架构}

\subsubsection*{2.1云-边-端三层架构}
\addcontentsline{toc}{subsubsection}{2.1云-边-端三层架构}

本文采用“云-边-端”三层协作架构。云端负责模型训练、全局调度和历史数据存储；边缘层部署Flink集群，负责实时流处理和本地决策；终端层负责数据采集和指令执行。

\subsubsection*{2.2 Flink流处理引擎}
\addcontentsline{toc}{subsubsection}{2.2 Flink流处理引擎}

在边缘节点上，我们部署了轻量化的Flink实例，配置了基于RocksDB的状态后端以支持大状态管理，并启用了检查点（Checkpoint）机制保障容错性。

\begin{quote}
{\fzht 状态后端配置} RocksDB状态后端能将状态数据持久化到本地磁盘，适合大状态场景。在边缘节点内存有限的条件下，这一特性尤为重要。
\end{quote}

\begin{quote}
{\fzht 检查点优化}采用增量检查点策略，仅保存自上次检查点以来的状态变更，大幅减少了检查点开销。
\end{quote}

\subsection*{三、动态资源调度策略}
\addcontentsline{toc}{subsection}{三、动态资源调度策略}

\subsubsection*{3.1问题建模}
\addcontentsline{toc}{subsubsection}{3.1问题建模}

边缘节点的计算资源有限且动态变化，因此资源调度问题可以建模为约束优化问题。设边缘节点集合为$\mathcal{N} = \{n_1, n_2, \dots, n_k\}$，任务集合为$\mathcal{T} = \{t_1, t_2, \dots, t_m\}$。

\begin{quote}
{\fzht 公式1}目标函数为最小化加权延迟与资源成本的组合：
\$$\min \sum_{i=1}^{m} \sum_{j=1}^{k} \bigl(w_1 \cdot x_{ij} \cdot L(t_i, n_j) + w_2 \cdot x_{ij} \cdot C(t_i, n_j)\bigr)$\$
\end{quote}

约束条件包括：

\begin{itemize}
  \item 每个任务必须分配到恰好一个节点：$\sum_{j} x_{ij} = 1, \forall i$
  \item 节点资源上限约束：$\sum_{i} x_{ij} \cdot r(t_i) \leq R(n_j), \forall j$
  \item 延迟SLA约束：$L(t_i, n_j) \leq L_{\max}, \forall (i,j): x_{ij}=1$
\end{itemize}

\subsubsection*{3.2算法设计}
\addcontentsline{toc}{subsubsection}{3.2算法设计}

采用改进的启发式贪心算法进行在线调度，如下所示。

\begin{quote}
{\fzht 算法1}在线动态资源调度算法

{\fzht 输入}：任务队列$\mathcal{Q}$，节点列表$\mathcal{N}$，SLA约束$L_{\max}$

{\fzht 输出}：任务-节点映射$\mathcal{M}$

``\inlinecode{
for each t in Q do
    C ← \{ n in N |满足资源约束\}
    if C is empty then
将t加入等待队列
        continue
    end if
    n* ← argmin\_\{n in C\} Score(t, n)
    if Score(t, n*) ≤ L\_max then
        M[t] ← n*
更新n*的资源状态
    else
将t升级到云端处理
    end if
end for
return M
}``
\end{quote}

\subsection*{四、实验评估}
\addcontentsline{toc}{subsection}{四、实验评估}

\subsubsection*{4.1实验环境}
\addcontentsline{toc}{subsubsection}{4.1实验环境}

实验使用3台边缘服务器（Intel Xeon E-2278G, 32GB RAM）和10台模拟终端设备。Flink集群配置为3个TaskManager，每个分配8GB堆内存。

\subsubsection*{4.2吞吐量对比}
\addcontentsline{toc}{subsubsection}{4.2吞吐量对比}

\begin{table}[htbp]
  \centering
  \small
  \begin{tabularx}{\textwidth}{>{\raggedright\arraybackslash}X >{\raggedright\arraybackslash}X >{\raggedright\arraybackslash}X}
    \toprule
    {\fzht 框架} & {\fzht 低负载（条/秒）} & {\fzht 高负载（条/秒）} \\
    \midrule
    Apache Storm & 12,500 & 38,200 \\
    Spark Streaming & 9,800 & 45,600 \\
    {\fzht 本文方法} & {\fzht 18,200} & {\fzht 60,800} \\
    \bottomrule
  \end{tabularx}
\end{table}

\subsubsection*{4.3延迟分析}
\addcontentsline{toc}{subsubsection}{4.3延迟分析}

在低负载（1000条/秒）条件下，本文方法的P99延迟为8.2ms，显著优于Storm的15.4ms和Spark Streaming的120ms（微批次导致）。在高负载（5000条/秒）条件下，本文方法通过动态调度策略仍将P99延迟控制在23.5ms以内。

\subsubsection*{4.4资源利用率}
\addcontentsline{toc}{subsubsection}{4.4资源利用率}

动态调度策略使CPU利用率从静态分配方案的62\%提升至78\%，内存利用率从55\%提升至72\%。

\subsection*{五、相关工作}
\addcontentsline{toc}{subsection}{五、相关工作}

边缘计算的概念最早由Shi等人系统阐述。在流处理方面，Carbone等人提出的Apache Flink以其真正的流处理语义成为业界标准。近年来，研究者们提出了多种边缘-云端协同的流处理方案。

与现有工作相比，本文的主要区别在于：

\begin{enumerate}
  \item 针对边缘资源受限场景设计了轻量化Flink配置
  \item 提出了在线动态资源调度算法，而非静态分配
  \item 在真实边缘服务器环境下进行了充分验证
\end{enumerate}

\subsection*{六、结论与展望}
\addcontentsline{toc}{subsection}{六、结论与展望}

本文提出了一种基于Apache Flink的边缘实时数据流处理框架，并设计了动态资源调度优化策略。实验结果表明，所提方法在吞吐量和资源利用率方面均有显著提升。

未来的研究方向包括：

\begin{itemize}
  \item 探索基于强化学习的智能调度策略
  \item 支持异构硬件（GPU/FPGA）的流处理加速
  \item 在更大规模边缘集群上进行验证
\end{itemize}

\begin{thebibliography}{99}
\bibitem{ref1} Shi W, Cao J, Zhang Q, et al. Edge computing: Vision and challenges[J]. IEEE Internet of Things Journal, 2016, 3(5): 637-646.
\bibitem{ref2} Carbone P, Katsifodimos A, Ewen S, et al. Apache Flink: Stream and batch processing in a single engine[J]. Bulletin of the IEEE Computer Society Technical Committee on Data Engineering, 2015, 38(4): 28-38.
\bibitem{ref3} Toshniwal A, Taneja S, Shukla A, et al. Storm@Twitter[C]. SIGMOD, 2014: 147-156.
\bibitem{ref4} Zaharia M, Xin R S, Wendell P, et al. Apache Spark: A unified engine for big data processing[J]. Communications of the ACM, 2016, 59(11): 56-65.
\end{thebibliography}

\end{document}